\documentclass[a4paper,12pt]{article}
\usepackage[T2A]{fontenc}
\usepackage[utf8]{inputenc}
\usepackage[russian]{babel}
\usepackage{amsmath,amssymb,mathtools}
\usepackage{geometry}
\geometry{left=23mm,right=21mm,top=21mm,bottom=23mm}
\usepackage{array,booktabs,tabularx}
\usepackage{xcolor,tikz}
\usepackage{hyperref}
\usepackage{fancyhdr}
\usepackage{enumitem}
\usepackage{microtype}
\definecolor{accent}{HTML}{007D82}
\definecolor{deepaccent}{HTML}{00565A}
\definecolor{ink}{HTML}{303A40}
\definecolor{err}{HTML}{AC3540}
\hypersetup{colorlinks=true,linkcolor=deepaccent,pdfauthor={Лёвин Артём Александрович},pdftitle={Ревью домашней работы — Кирилл Зиновьев}}
\setlength{\parindent}{0pt}
\setlength{\parskip}{7pt plus 1pt}
\renewcommand{\arraystretch}{1.35}
\pagestyle{fancy}
\fancyhf{}
\fancyhead[L]{\small\color{ink}Линейная функция}
\fancyhead[R]{\small\color{ink}Кирилл Зиновьев}
\fancyfoot[C]{\small\thepage}
\renewcommand{\headrulewidth}{0.3pt}
\setlength{\headheight}{15pt}
\newcommand{\topic}[1]{\textcolor{deepaccent}{\textbf{#1}}}
\newcommand{\flagerr}{\textcolor{err}{\textbf{Ошибка:}}}
\newcommand{\flagmiss}{\textcolor{err}{\textbf{Не представлено:}}}
\newcommand{\idea}{\textcolor{deepaccent}{\textbf{Ключевая идея:}}}
\newcommand{\finalidea}{\textcolor{accent}{\textbf{Итог:}}}
\newcommand{\taskhead}[2]{\phantomsection\label{task:#1}{\Large\color{deepaccent}\bfseries Задание № #1}\par\medskip{\large #2}\par\smallskip\hrule\medskip}
\begin{document}
\thispagestyle{empty}
\begin{center}
\vspace*{24mm}
{\large\color{deepaccent} МАТЕМАТИКА}\par
\vspace{23mm}
{\LARGE Ревью домашней работы}\par
\vspace{15mm}
{\Large\color{accent} Линейная функция и пересечение прямых}\par
\vspace{9mm}
{\normalsize Угловой коэффициент, подстановка точки, оси координат}\par
\vspace{42mm}
{\large Кирилл Зиновьев}\par
\vspace{13mm}
{17 ноября 2026 г.}\par
\vfill
{\small Лёвин Артём Александрович}\par
{\small эксклюзивно для Кирилла Зиновьева}\par
\vspace{12mm}
\begin{tikzpicture}
\draw[accent,line width=0.7pt] (0,0) .. controls (1.2,0.55) and (2.0,-0.35) .. (3.15,0)
.. controls (4.15,0.4) and (5.2,-0.45) .. (6.3,0);
\end{tikzpicture}
\end{center}
\clearpage
\section*{Сводная таблица}
В работе проверены все десять заданий с учётом дополнительной фотографии решений № 6--8. Подробного разбора требует только задание № 10. Остальные девять заданий выполнены верно и в таблицу не включены.

\smallskip
{\small
\begin{tabularx}{\linewidth}{@{}p{16mm}p{29mm}>{\raggedright\arraybackslash}X p{25mm}p{25mm}@{}}
\toprule
№ & Тема & Суть ошибки & Ваш ответ & Верный ответ\\
\midrule
\hyperref[task:10]{10} & Три прямые & При вычислении $y$ пропущен множитель $2$ & $k=\frac{10}{3}$ & $k=\frac{7}{3}$\\
\bottomrule
\end{tabularx}
}
\vspace{7mm}
\topic{На что обратить внимание.} После нахождения абсциссы точки пересечения следует подставить её в обе исходные формулы и сравнить полученные ординаты. Этот короткий шаг помогает обнаружить арифметическую ошибку до вычисления неизвестного коэффициента.
\clearpage
\thispagestyle{empty}
\vspace*{0.33\textheight}
\begin{center}
{\Huge\color{deepaccent}Разбор ошибок}\par\vspace{8mm}
{\large Задание № 10: точная подстановка координат}\par
\end{center}
\clearpage
\begingroup\small\setlength{\parskip}{3pt}\setlength{\abovedisplayskip}{5pt}\setlength{\belowdisplayskip}{5pt}\setlength{\abovedisplayshortskip}{3pt}\setlength{\belowdisplayshortskip}{3pt}
\taskhead{10}{Поиск общего пересечения трёх прямых}
\textbf{Текст задания.} Найдите $k$, если прямая $y=kx-3$ проходит через точку пересечения прямых $y=10-2x$ и $y=x+1$.

\textbf{Ваше решение.} В тетради верно найдено $x=3$, после чего записано:
\[
\begin{aligned}
10-2x=x+1&\ \Longrightarrow\ 9=3x\ \Longrightarrow\ x=3,\\
\textcolor{err}{10-3}&=\textcolor{err}{7},\qquad \textcolor{err}{(x;y)=(3;7)},\\
7&=3k-3\ \Longrightarrow\ k=\frac{10}{3}.
\end{aligned}
\]
\textbf{Ваш ответ:} $k=\frac{10}{3}$.

\flagerr\ При подстановке $x=3$ в выражение $10-2x$ пропущен множитель $2$.

\textbf{Где именно возникает ошибка.} Последний полностью правильный шаг --- $x=3$. Следующая записанная строка $10-3=7$ уже неверно использует формулу прямой. Правильная подстановка:
\[
10-2\cdot3=10-6=4.
\]
\textbf{Почему это неверно.} В формуле $y=10-2x$ абсцисса умножается на $-2$. При значении $x=3$ нужно вычислить $2\cdot3=6$, затем вычесть $6$ из $10$. Если заменить $2x$ на $x$, получится другая функция $y=10-x$, которая здесь не дана.

\idea\ Точка пересечения двух прямых должна удовлетворять \emph{обоим} их уравнениям. Сначала находят $x$, затем --- $y$, проверяют координаты в обеих формулах и только после этого определяют неизвестный коэффициент третьей прямой.

\textbf{Исправление от последнего верного шага.} Сохраняем найденное $x=3$. Тогда
\[
y=10-2\cdot3=4,
\qquad\text{следовательно,}\quad (x;y)=(3;4).
\]
\textbf{Полное правильное решение.}
\begin{enumerate}[leftmargin=*,itemsep=4pt]
\item При пересечении ординаты равны:
\[
10-2x=x+1.
\]
\item Переносим все слагаемые с $x$ вправо, числа --- влево:
\[
10-1=x+2x,\qquad 9=3x,\qquad x=3.
\]
\item Находим ординату по любой из исходных функций:
\[
y=x+1=3+1=4.
\]
\item Подставляем найденную точку $(3;4)$ в третью прямую $y=kx-3$:
\[
4=3k-3,\qquad 7=3k,\qquad k=\frac73.
\]
\end{enumerate}
\textbf{Проверка.} При $x=3$ обе исходные прямые дают одно значение:
\[
10-2\cdot3=4,\qquad 3+1=4.
\]
Найденный коэффициент также подходит:
\[
\frac73\cdot3-3=7-3=4.
\]
Все три прямые проходят через $(3;4)$.

\finalidea\ $\boxed{k=\frac73}$.

\topic{Задача на закрепление.} Найдите $m$, если прямая $y=mx-4$ проходит через точку пересечения прямых $y=11-2x$ и $y=2x-1$.
\endgroup
\end{document}